% !TeX program = xelatex
% !TEX encoding = UTF-8

\documentclass[a4paper,12pt]{article}
\usepackage[margin=2.2cm,headheight=16pt]{geometry}
\usepackage{amsmath,amssymb}
\usepackage{booktabs,longtable,tabularx,array}
\usepackage[AutoFakeBold,AutoFakeSlant,CJKmath=true]{xeCJK}
\usepackage{fancyhdr}
\usepackage{xcolor}
\usepackage{enumitem}
\usepackage{ragged2e}
\usepackage{setspace}
\usepackage[colorlinks=true,linkcolor=black,urlcolor=blue]{hyperref}

\setCJKmainfont{黑體-繁}
\setCJKsansfont{黑體-繁}
\setCJKmonofont{黑體-繁}

\definecolor{NtmcBlue}{HTML}{1F4E79}
\hypersetup{
  pdftitle={新北捷運人員月排班問題：數學模型與 CP-SAT 解法},
  pdfauthor={吳亞倫},
  pdfsubject={三鶯站務 M、環狀站務 YM、三鶯檢修 T 與環狀檢修 YT 月排班模型}
}

\pagestyle{fancy}
\fancyhf{}
\fancyhead[L]{新北捷運人員月排班問題}
\fancyhead[R]{數學模型與 CP-SAT 解法}
\fancyfoot[C]{\thepage}
\renewcommand{\headrulewidth}{0.4pt}
\setstretch{1.16}
\setlength{\parindent}{2em}
\setlength{\parskip}{3pt plus 1pt minus 1pt}
\setlist{itemsep=2pt,topsep=4pt}
\setlength{\extrarowheight}{2pt}
\renewcommand{\arraystretch}{1.2}
\renewcommand\tabularxcolumn[1]{>{\RaggedRight\arraybackslash}p{#1}}
\newcolumntype{Y}{>{\centering\arraybackslash}X}
\renewcommand{\abstractname}{摘要}
\renewcommand{\contentsname}{目錄}

\begin{document}
\begin{titlepage}
  \centering
  \vspace*{3.2cm}
  {\Large\bfseries 新北捷運公司人員自動排班系統\par}
  \vspace{1.2cm}
  {\LARGE\bfseries\color{NtmcBlue} 新北捷運人員月排班問題\par}
  \vspace{0.4cm}
  {\Large 數學模型與 CP-SAT 解法\par}
  \vfill
  \begin{tabular}{rl}
    撰寫者：& 吳亞倫（1508S0342）\\
    單位：& 企劃處計畫室規劃課\\
    版本日期：& 2026 年 8 月 26 日
  \end{tabular}
  \vfill
\end{titlepage}

\begin{abstract}
  本文研究都市軌道運輸站務與檢修人員之月排班問題，並分別建立適用於三鶯線與環狀線
  站務單位（M、YM）及檢修單位（T、YT）的限制最佳化模型。模型以完整前月班表、目標月
  人力需求、人員屬性、指定休假、固定勤務及八週休假週期為已知資料，決定每位人員每日
  的工作地點、班別與休假種類。可行域同時刻畫每日唯一指派、站點覆蓋、最少十一小時
  休息、連續七日一般休假及八週休假額度等條件；班組人力品質與其他偏好則轉化為非負
  整數違反量，並以
  字典序目標維持不同管理目的間的優先關係。求解採用 CP-SAT，透過布林變數、整數變數、
  重述限制與分段線性化表達離散排班邏輯。為妥善處理月界，規劃範圍向次月延伸七日，並
  由歷史班表初始化跨月狀態。除最優化主班表外，本文亦定義提示解、有限時間搜尋及具最小
  漢明距離之替代候選生成方法。所建模型兼顧法規可行性、營運覆蓋、工作型態與群組內公平，
  且各項目標皆可分解為明確的違反量與權重，因而具備可驗證性與可解釋性。
\end{abstract}

\noindent\textbf{關鍵詞：}人員排班、限制規劃、CP-SAT、字典序最佳化、休假公平、候選解

\tableofcontents
\clearpage

\section{問題背景與研究範圍}

人員排班的核心是在法規、營運需求與人員偏好同時存在時，於有限規劃期間內建立
逐人逐日的可行指派。硬性條件之間具有顯著的時間耦合：某日班別會影響隔日可選班別，
一般休假會改變連續日數狀態，八週額度又跨越單一月份。另一方面，指定休假、工作區段
長度、班別穩定性及公平性通常無法同時達到各自理想值，因此須以明確的優先結構描述
其取捨，而不能只使用一個未經尺度校準的總和目標。

本研究考慮四個排班單位：三鶯站務 M、環狀站務 YM、三鶯檢修 T 與環狀檢修 YT。
M 與 YM 具有相同的數學結構，差異由站點集合、站群、班位需求及班別時間參數給定；
T 與 YT 亦共享同一檢修模型結構，而人員、需求與班組資料彼此獨立。站務模型的決策
包含工作站點與班別，並允許以外援補足特定站點；檢修模型則以月班組為基礎，兼顧
正常出勤、專業組別與能力配置。兩類模型共同決定一般休假 \(R\)、國定假日休假
\(R1\)、自選休假 \(R休\) 及固定勤務 \(X\)。

規劃期間由目標月份 \(\mathcal D^{\mathrm{tar}}\) 與其後七日
\(\mathcal D^+\) 組成。延伸區間使月底決策仍須滿足休息間隔、連續日數及八週額度的
可延續性；僅目標月指派列入正式班表與月內統計。前月完整班表則提供規劃起點以前的
狀態，使跨月限制不因任意截斷時間軸而失真。全體時間均以台北時間（UTC+8）表示，
跨午夜勤務歸屬其開始日期。

\subsection{輸入資料與邊界}

模型輸入視為求解開始時固定的資料快照，包含目標月份、人員集合與屬性、日別需求、
完整前月實際班表、連續的 56 日休假區間、指定休假、固定指派，以及固定勤務 \(X\)
的實際起訖時間。對規劃起點前已在職的人員，前月每一日在職日皆須有唯一歷史狀態；
目標月內到職者則以到職日作為其時間序列起點。此完備性假設確保所有跨月狀態均可由
觀測資料唯一初始化。

月班表的日格可表示正常班、\(R\)、\(R1\)、\(R休\)、\(R^*\) 或
\(X[\mathrm{HH{:}mm-HH{:}mm}\mid\text{註記}]\)。\(R^*\) 表示待模型決定底層休假種類的
指定休假；若申請獲得滿足，實際狀態必為 \(R\)、\(R1\) 或 \(R休\) 之一，並可在輸出中
同時保留申請標記。\(R休\) 只可用於目標月的 \(R^*\)，每人給定數量為使用上限。
\(X\) 的文字註記只具識別作用，
不參與任何限制、目標函數或可行性判定。

模型輸出為目標月份內的完整指派，以及各軟性條件的違反量、權重與加權分數。
一個候選班表稱為\emph{可行}，若其滿足全部硬性限制；稱為\emph{字典序最優}，若不存在
另一可行班表能在第一個相異的目標群組取得更小值。此區分在有限時間求解中特別重要：
已找到的 incumbent 可保證可行，但僅在各階段均取得最優性證明時，方可宣稱整體字典序最優。

\section{建模與求解方法}

\subsection{CP-SAT 與字典序架構}

令決策向量 \(\mathbf x\) 收集所有布林與整數變數，\(\mathcal{F}\) 表示全部硬性
限制所形成的有限可行域。每一項軟性條件定義一個
非負違反量
\[
  \Phi_j(\mathbf{x})\in\mathbb{Z}_{\geq0}.
\]
同一優先群組內的規則以加權和組成
\[
  J_p(\mathbf{x})
  =
  \sum_{j\in\mathcal{J}_p}\omega_j\Phi_j(\mathbf{x}),
\]
不同群組之間採字典序最佳化：
\[
  \operatorname{lexmin}_{\mathbf{x}\in\mathcal{F}}
  \left(J_1(\mathbf{x}),J_2(\mathbf{x}),\ldots,J_P(\mathbf{x})\right).
\]
亦即先最小化 \(J_1\)，固定其最優值後再最小化 \(J_2\)，依此類推。群組權重
只表示同一群組內不同違反量的交換比例，不能跨群組抵銷較高優先目標的惡化。

字典序問題以逐階段 CP-SAT 求解實現。令第 \(p\) 階段所得最優值為 \(J_p^*\)，
則下一階段的可行域遞迴定義為
\[
  \mathcal F_1=\mathcal F,
  \qquad
  \mathcal F_{p+1}
  =\{\mathbf x\in\mathcal F_p:J_p(\mathbf x)=J_p^*\}.
\]
第 \(p\) 階段求解
\[
  J_p^*=\min_{\mathbf x\in\mathcal F_p}J_p(\mathbf x).
\]
由此可知，一旦某階段值以等式固定，後續階段無法犧牲該值；若所有階段皆完成最優性
證明，最後所得解即為原問題的字典序最優解。

CP-SAT 要求目標與限制以整數線性式、布林邏輯或其全域限制表示。本文中的蘊含關係
\(b=1\Rightarrow a^\top\mathbf x\leq c\) 以條件式限制重述；絕對值
\(u=|v|\) 以最大值等式 \(u=\max(v,-v)\) 表示；分段函數則以有限整數定義域上的
元素查表限制編碼。平方偏差的自變數皆具有有限整數範圍，故亦可預先列舉函數值，
而不引入連續變數或近似誤差。所有係數經整數尺度化，因此模型保持為純整數限制最佳化問題。

\section{站務 M／YM 共用模型}

站務模型以參數化方式同時描述 M 與 YM。兩者使用相同的變數、可行域與目標函數，
並分別代入所屬路線的站點集合、站群劃分、班位上下限、外援等級及班別時間。
因此以下以單一符號系統推導，求解特定單位時再以其資料實例化各集合與參數。

\subsection{輸入與輸出}

M 模型的輸入包括：
\begin{itemize}
  \item 目標月份與固定七日延伸區間；
  \item M 人員及其所屬車站；
  \item 指定休假 \(R^*\)、每人目標月 \(R休\) 使用上限、固定勤務 \(X\) 與其他固定指派；
  \item 八週週期的一般 \(R\) 與 \(R1\) 額度；
  \item 依每人在職日期推導的目標月份一般 \(R\) 與 \(R1\) 基準數量；
  \item 完整前月實際班表，用以重建所有跨月邊界狀態；
  \item 建模開始前各人員、各八週週期已使用的 R 與 R1 數量；
  \item 國定假日集合。星期六與星期日由日期本身決定。
\end{itemize}

模型輸出每位 M 人員在目標月份內的實際狀態。正常班需同時指出車站與班別；
滿足 \(R^*\) 時，輸出仍保留申請標記，並同時記錄其底層實際使用的是 \(R\)、
\(R1\) 或 \(R休\)。延伸日期只屬於模型內部決策，不列入正式輸出。

\subsection{集合與索引}

\begin{description}
  \item[\(\mathcal{E}^{M}\)] M 人員集合，索引為 \(e\)。
  \item[\(\mathcal{D}^{\mathrm{tar}}\)] 目標月份內的日期集合，索引為 \(d\)。
  \item[\(\mathcal{D}^{+}\)] 目標月份之後固定七日的延伸日期集合。
  \item[\(\mathcal{D}\)] 完整建模日期集合，
    \(\mathcal{D}=\mathcal{D}^{\mathrm{tar}}\cup\mathcal{D}^{+}\)。
  \item[\(\mathcal{L}\)] 車站集合，索引為 \(l\)。
  \item[\(\mathcal{S}\)] 正常班別集合，
    \(\mathcal{S}=\{\mathrm{E},\mathrm{A},\mathrm{N}\}\)，分別代表早、午、夜班。
  \item[\(\mathcal{C}\)] 八週休假週期集合，索引為 \(c\)。
  \item[\(\mathcal{G}\)] 車站群組集合。
  \item[\(\mathcal{D}^{wd}\)] 平日集合，即星期一至星期五且不是國定假日的日期。
  \item[\(\mathcal{D}^{hol}\)] 假日集合，即星期六、星期日或國定假日的日期。
  \item[\(\mathcal{L}^{ext}\)] 允許使用外派人力的車站集合。
\end{description}

M 車站集合固定為 \(\mathcal L_M=\{\mathrm{LB01},\ldots,\mathrm{LB12}\}\)，
YM 固定為 \(\mathcal L_{YM}=\{\mathrm{Y06},\ldots,\mathrm{Y19}\}\)；求解時 \(\mathcal L\)
取所屬工作區的固定集合。
群組函數、班位上下限與外援等級由目標月份快照提供，群組非空但不固定站數。
令 \(\mathcal L^{disc}\) 與 \(\mathcal L^{allow}\) 分別為「盡量不要」及
「允許」外援的車站，且
\(\mathcal L^{ext}=\mathcal L^{disc}\cup\mathcal L^{allow}\)。

\subsection{參數}

對每位人員 \(e\)，令 \(h_e\in\mathcal{L}\) 為所屬車站，\(g(l)\) 為車站
\(l\) 的群組，並定義合法工作車站
\[
  \mathcal{A}_e
  =
  \{l\in\mathcal{L}:g(l)=g(h_e)\}.
\]

令 \(b^-_{ls},b^+_{ls}\in\mathbb{Z}_{\geq0}\) 為月份快照所定的班位人數
下限與上限，且 \(b^-_{ls}\le b^+_{ls}\)。同一設定適用目標月與月底延伸日。

M 的班別時間為：早班 \(06{:}30\)--\(14{:}30\)、午班
\(14{:}20\)--\(22{:}20\)、夜班 \(22{:}00\)--翌日 \(07{:}00\)。夜班歸屬
其開始日期。

令 \(p_{ed}\in\{0,1\}\) 表示人員 \(e\) 是否在日期 \(d\) 提出指定休假
\(R^*\)。\(R^*\) 不是獨立休假類型；若申請被滿足，該日仍須實際使用
\(R\)、\(R1\) 或 \(R休\)。令 \(k_e\in\mathbb{Z}_{\geq0}\) 為人員 \(e\)
在目標月份可使用的 \(R休\) 上限，取自需求列的 \texttt{當月指定R休}；
歷史與候選列的同名欄位則記錄日格實際 \(R休\) 數。令 \(z_{ed}\in\{0,1\}\) 表示該日是否具有固定勤務 \(X\)，
且每筆 \(X\) 具有固定的實際開始與結束時間。

對每個八週週期 \(c\)，令 \(G_c\) 為一般 \(R\) 額度，令 \(H_c\) 為 \(R1\)
額度；本模型採 \(G_c=16\)。令
\(\overline{R}_{ec}\)、\(\overline{R1}_{ec}\) 分別為建模區間開始前，週期
\(c\) 已使用的一般 \(R\) 與 \(R1\) 數量。令月份設定中的共同基準為
\(N^R,N^{R1}\)，令 \(u^R_e,u^{R1}_e\) 分別為人員 \(e\) 到職日前、但仍在
目標月份內的週六日與國定假日數，並令
\[
  n_e=\max\{0,N^R-u^R_e\},
  \qquad
  m_e=\max\{0,N^{R1}-u^{R1}_e\}.
\]
\(N^R\) 取當月週六日數，\(N^{R1}\) 取當月國定假日數；兩者亦可作為外生參數，
但必須位於 56 日硬額度所允許的共同範圍內。這兩者只影響軟目標，不改變八週精確額度。

完整前月歷史班表用來重建法規連續日數、尚未結束的實際工作區段、
該工作區段最近一次正常班型與是否已混合班型。八週 R/R1 已使用數量由獨立
輸入提供，不從七日歷史推算；M／YM／T／YT 既有人員若在 Demand 明確
提供月初 R/R1，該值優先於前月月底累計，僅在留空時由前月歷史推導；
新進人員仍依到職日前日期折抵。公平性不使用歷史累積值。

\subsection{決策變數與衍生狀態}

主要決策變數為
\[
  x^{M}_{edls}\in\{0,1\},
\]
其中 \(x^{M}_{edls}=1\) 表示人員 \(e\) 在日期 \(d\) 於車站 \(l\) 上班別
\(s\)。休假變數為
\[
  r^{M}_{ed},\ r^{M,1}_{ed},\ r^{M,L}_{ed}\in\{0,1\},
\]
分別表示一般 \(R\)、\(R1\) 與 \(R休\)。外派班位變數為
\[
  a_{dls}\in\mathbb{Z}_{\geq0}.
\]

定義衍生變數
\[
  y^{M}_{eds}=\sum_{l\in\mathcal{A}_e}x^{M}_{edls},
  \qquad
  w^{M}_{ed}=\sum_{s\in\mathcal{S}}y^{M}_{eds}+z_{ed},
  \qquad
  \rho^{M}_{ed}=r^{M}_{ed}+r^{M,1}_{ed}+r^{M,L}_{ed},
\]
其中 \(w^{M}_{ed}=1\) 表示實際工作日，\(\rho^{M}_{ed}=1\) 表示實際休假日。
另令
\[
  u_{ed}
  =
  \sum_{s\in\mathcal{S}}
  \sum_{\substack{l\in\mathcal{A}_e\\l\neq h_e}}x^{M}_{edls}
\]
表示該日是否跨站支援。

為避免混淆，下列狀態變數具有不同用途：

\begin{center}
\small
\begin{tabularx}{\textwidth}{p{2.6cm}X X}
  \toprule
  狀態 & 計數內容與中斷條件 & 用途 \\
  \midrule
  \(q^{law}_{ed}\) & 自最近一次一般 \(R\) 後經過的排班日數；只有一般 \(R\) 歸零，\(R1\) 與 \(R休\) 仍增加 & 法規性最多連續六日 \\
  \(c^{work}_{ed}\) & 連續實際工作日數；正常班與 \(X\) 增加，\(R\)、\(R1\) 或 \(R休\) 歸零 & 工作與休假區段長度 \\
  \((\beta_{ed},\mu_{ed})\) & 實際工作區段內最近正常班型及是否已出現多種班型；休假歸零，\(X\) 保留狀態 & 工作區段班型一致性 \\
  \(\gamma_{ed}\) & 最近一次實際休假後出現的最近正常班別；\(R\)、\(R1\)、\(R休\) 清除，\(X\) 略過 & 換班前是否有休假 \\
  \(\eta_{ed}\) & 全序列中最近一次有效正常班別；\(R\)、\(R1\)、\(R休\)、\(X\) 都略過 & 班別輪轉方向 \\
  \bottomrule
\end{tabularx}
\normalsize
\end{center}

另定義三日模式指示變數
\(\pi^{NE}_{ed},\pi^{NA}_{ed}\in\{0,1\}\)，分別表示夜--休--早與夜--休--午。

\subsection{硬性限制}

\paragraph{每日唯一指派}
每位人員每天恰好具有一個實際狀態：
\[
  \sum_{l\in\mathcal{A}_e}\sum_{s\in\mathcal{S}}x^{M}_{edls}
  +r^{M}_{ed}+r^{M,1}_{ed}+r^{M,L}_{ed}+z_{ed}=1,
  \qquad
  \forall e\in\mathcal{E}^{M},\ d\in\mathcal{D}.
\]
固定正常班、固定 \(R\)、固定 \(R1\)、固定 \(R休\) 與固定 \(X\) 直接固定對應變數。
\(R^*\) 只表示休假偏好，不固定底層休假類型。

\paragraph{每月 R休 上限}
\[
  r^{M,L}_{ed}\leq p_{ed},
  \qquad
  \sum_{d\in\mathcal{D}^{\mathrm{tar}}}r^{M,L}_{ed}\leq k_e.
\]
因此 \(R休\) 只能使用目標月的 \(R^*\)，延伸日不安排 \(R休\)，亦不形成跨月額度。

\paragraph{合法工作站點}
對所有 \(l\notin\mathcal{A}_e\)，要求
\[
  x^{M}_{edls}=0.
\]
因此人員只能在本站或同群組車站工作。

\paragraph{班位覆蓋}
內部人員與外援合計必須落在月份設定的閉區間：
\[
  b^-_{ls}\le
  \sum_{e\in\mathcal{E}^{M}}x^{M}_{edls}+a_{dls}
  \le b^+_{ls},
  \qquad
  \forall d\in\mathcal{D},\ l\in\mathcal{L},\ s\in\mathcal{S}.
\]

\paragraph{外援站點限制}
對所有 \(l\notin\mathcal{L}^{ext}\)，要求
\[
  a_{dls}=0.
\]
外派只代表外部人力補足一個班位，不對應任何 M 人員列。

\paragraph{最少十一小時休息}
令 \(\mathcal{I}^{M}\) 為所有實際休息時間少於十一小時的正常班組合。對任一
\(((d,s),(d',s'))\in\mathcal{I}^{M}\)，要求
\[
  y^{M}_{eds}+y^{M}_{ed's'}\leq1,
  \qquad
  \forall e\in\mathcal{E}^{M}.
\]
若正常班與固定 \(X\) 的實際區間不足十一小時，該正常班變數固定為 0；兩筆固定
\(X\) 若彼此重疊或間隔不足十一小時，輸入本身不可行。此規則只比較實際工作
結束與下一次實際工作開始，與班別名稱是否相同無關。

\paragraph{連續七日至少一日一般 R}
此規則是法規性連續日數限制。對每日 \(d\)，要求
\[
  r^{M}_{ed}=1
  \ \Longrightarrow\ 
  q^{law}_{ed}=0,
\]
以及
\[
  r^{M}_{ed}=0
  \ \Longrightarrow\ 
  q^{law}_{ed}=q^{law}_{e,d-1}+1.
\]
模型限制
\[
  0\leq q^{law}_{ed}\leq6.
\]
因此正常班、\(X\)、\(R1\) 與 \(R休\) 都會使計數增加；只有一般 \(R\) 會歸零。
例如下列序列的計數會到 7，故不合法：
\[
  (\mathrm{E},\mathrm{E},\mathrm{E},R1,\mathrm{A},\mathrm{A},\mathrm{A})
  \longmapsto (1,2,3,4,5,6,7).
\]
若 \(R^*\) 以 \(R\) 滿足則歸零，以 \(R1\) 滿足則不歸零。

\paragraph{八週休假額度}
對每個週期 \(c\)，令
\[
  \mathcal{D}_c=c\cap\mathcal{D},
  \qquad
  F_c=\left|\{d\in c:d>\max\mathcal{D}\}\right|,
\]
其中 \(F_c\) 是模型區間結束後、週期內尚未建模的日期數。

若模型已涵蓋週期 \(c\) 的結束日，則
\[
  \overline{R}_{ec}
  +\sum_{d\in\mathcal{D}_c}r^{M}_{ed}
  =G_c,
  \qquad
  \overline{R1}_{ec}
  +\sum_{d\in\mathcal{D}_c}r^{M,1}_{ed}
  =H_c.
\]

若模型尚未涵蓋週期結束日，則兩種額度均不得超用：
\[
  \overline{R}_{ec}+\sum_{d\in\mathcal{D}_c}r^{M}_{ed}\leq G_c,
  \qquad
  \overline{R1}_{ec}+\sum_{d\in\mathcal{D}_c}r^{M,1}_{ed}\leq H_c,
\]
並且剩餘日期必須分別足以補足兩種額度：
\[
  \overline{R}_{ec}+\sum_{d\in\mathcal{D}_c}r^{M}_{ed}+F_c\geq G_c,
\]
\[
  \overline{R1}_{ec}+\sum_{d\in\mathcal{D}_c}r^{M,1}_{ed}+F_c\geq H_c.
\]
由於同一日不能同時使用 \(R\) 與 \(R1\)，另需滿足聯合可補足性：
\[
  \overline{R}_{ec}+\overline{R1}_{ec}
  +\sum_{d\in\mathcal{D}_c}\left(r^{M}_{ed}+r^{M,1}_{ed}\right)
  +F_c
  \geq G_c+H_c.
\]
\(R1\) 可安排在週期內任意日期，不必安排於國定假日當日。\(R休\) 不計入上述任一額度。

\paragraph{固定指派與歷史邊界}
所有固定指派均不得由模型改寫。歷史班表決定建模區間第一日前的
\(q^{law}\)、\(c^{work}\)、\((\beta,\mu)\)、\(\gamma\)、\(\eta\) 及週期
累積值，這些值作為第一日狀態轉移的固定前置狀態。

\subsection{軟性目標與權重}

\subsubsection{\texorpdfstring{$J^M_1$}{J1}：指定休假}

\paragraph{指定休假滿足}
只要指定休假日實際使用 \(R\)、\(R1\) 或 \(R休\)，即視為申請被滿足。未滿足數為
\[
  \Phi^{M,R^*}
  =
  \sum_{\substack{e\in\mathcal{E}^{M},\ d\in\mathcal{D}^{\mathrm{tar}}\\p_{ed}=1}}
  \left(1-r^{M}_{ed}-r^{M,1}_{ed}-r^{M,L}_{ed}\right).
\]

\paragraph{未使用 R休 額度}
未使用的 \(R休\) 上限違反量為
\[
  \Phi^{M,unusedL}
  =\sum_{e\in\mathcal{E}^{M}}
  \left(k_e-\sum_{d\in\mathcal{D}^{\mathrm{tar}}}r^{M,L}_{ed}\right).
\]

\subsubsection{\texorpdfstring{$J^M_4$}{J4}：綜合排班品質}

\paragraph{外援人力}
令「允許」與「盡量不要」車站的目標月外援總人次分別為
\[
  A^{M,allow}
  =
  \sum_{d\in\mathcal{D}^{\mathrm{tar}}}
  \sum_{l\in\mathcal L^{allow}}
  \sum_{s\in\mathcal{S}}a_{dls}.
\]
\[
  A^{M,disc}
  =
  \sum_{d\in\mathcal{D}^{\mathrm{tar}}}
  \sum_{l\in\mathcal L^{disc}}
  \sum_{s\in\mathcal{S}}a_{dls}.
\]
「允許」外援共用前 70 人次免罰額度；「盡量不要」從第一人次起計分，故
\[
  \Phi^{M,ext}=\max\{0,A^{M,allow}-70\}+A^{M,disc}.
\]

\paragraph{每月一般 R 分布}
令
\[
  Q^{M,R}_e=\sum_{d\in\mathcal{D}^{\mathrm{tar}}}r^{M}_{ed}.
\]
定義
\[
  f_t(q)=|q-t|^2.
\]
則一般 \(R\) 的月內分布違反量為
\[
  \Phi^{M,month,R}
  =
  \sum_{e\in\mathcal{E}^{M}}f_{n_e}\!\left(Q^{M,R}_e\right),
\]
偏離基準 0、1、2、3 日時，懲罰依序為 0、1、4、9。
此為月內分布偏好，不取代八週休假額度限制；\(R休\) 不計入一般 \(R\) 數量。

\paragraph{八週累積 R1 餘額}
對每個與目標月相交的 56 日區間 \(I\)，令
\(p^{M,1}_{eI}\) 為目標月內可排班日前已使用或折抵的 \(R1\) 數。將 \(m_e\)
依各相交區間在可排日期內的國定假日數比例分配為 \(a_{eI}\)；若所有區間皆無
國定假日，改依相交日數比例分配，整數餘數依區間起日順序放入前段。並定義
\[
  A^{M,1}_{eI}=p^{M,1}_{eI}
  +\sum_{d\in\mathcal D^{\mathrm{tar}}\cap I}r^{M,1}_{ed},
  \qquad
  T^{M,1}_{eI}=|\{h\in\mathcal H_I:h<\max(d_{start,e},d_{month})\}|+a_{eI}.
\]
累積餘額的懲罰函數為
\[
  g(a,h)=\max\{0,a-h\}^2+\max\{0,h-a-1\}^2,
\]
故
\[
  \Phi^{M,balance,R1}=\sum_{e\in\mathcal E^M}\sum_{I}g(A^{M,1}_{eI},T^{M,1}_{eI}).
\]
累積剛好或暫欠一日皆不懲罰；提前多休與超過一日的欠額按平方計分。餘額沿同一
區間跨月承接，區間結束時仍由硬性限制強制精確額度。

\paragraph{連續工作區段長度}
此違反量控制實際工作與休假的分段品質，與只有一般 \(R\) 才重置的法規連續日數不同。
實際工作區段是連續曆日皆為正常班或 \(X\) 的最大區段；一般 \(R\)、
\(R1\) 與 \(R休\) 都會結束區段。定義區段長度懲罰
\[
  D^M(L)
  =
  \begin{cases}
    4, & L=1,\\
    0, & L=2,\\
    0, & L=3,\\
    0, & L=4,\\
    0, & L=5,\\
    2(L-4), & L\geq6.
  \end{cases}
\]
令 \(\mathcal{B}^{work}_e\) 為人員 \(e\) 的最大實際工作區段集合，並承接月初
未結束區段。只對結束日在目標月份內的區段計分：
\[
  \Phi^{M,streak}
  =
  \sum_{e\in\mathcal{E}^{M}}
  \sum_{B\in\mathcal{B}^{work}_e:\ \max B\in\mathcal{D}^{\mathrm{tar}}}
  D^M(|B|).
\]

\paragraph{工作區段班型一致性}
此違反量沿用實際工作區段 \(\mathcal{B}^{work}_e\) 與跨月承接。
對每個區段 \(B\)，令其中出現的正常班型集合為
\[
  \mathcal S_B
  =\{s\in\mathcal S:\exists d\in B,\ y^M_{eds}=1\}.
\]
若 \(|\mathcal S_B|\geq2\)，令 \(\mu_B=1\)，否則 \(\mu_B=0\)。\(X\) 保留在
工作區段內，但不加入 \(\mathcal S_B\)。只對結束日在目標月份內的區段計分：
\[
  \Phi^{M,mixed}
  =\sum_{e\in\mathcal E^M}
  \sum_{B\in\mathcal B^{work}_e:\ \max B\in\mathcal D^{\mathrm{tar}}}\mu_B.
\]

\paragraph{夜班--休假--早班型態}
對所有與目標月份相交的三日視窗（起日自月初前兩日至月底），令
\[
  \pi^{NE}_{ed}=1
  \Longleftrightarrow
  y^{M}_{ed\mathrm{N}}=1,
  \ \rho^{M}_{e,d+1}=1,
  \ y^{M}_{e,d+2,\mathrm{E}}=1.
\]
總違反量為
\[
  \Phi^{M,NE}=\sum_{e,d}\pi^{NE}_{ed}.
\]
此規則特別避免夜班後只休一日即接早班。

\paragraph{夜班--休假--午班型態}
同理，令
\[
  \pi^{NA}_{ed}=1
  \Longleftrightarrow
  y^{M}_{ed\mathrm{N}}=1,
  \ \rho^{M}_{e,d+1}=1,
  \ y^{M}_{e,d+2,\mathrm{A}}=1,
\]
並令
\[
  \Phi^{M,NA}=\sum_{e,d}\pi^{NA}_{ed}.
\]

\paragraph{未休假直接換班}
此規則判斷班別改變前是否出現實際休假。\(R\)、\(R1\) 或 \(R休\) 會清除狀態，\(X\)
只被略過。若當日正常班別為 \(s\)，且
\(\gamma_{e,d-1}\notin\{0,s\}\)，則產生一次違反。總違反量記為
\[
  \Phi^{M,restswitch}.
\]
因此 \(\mathrm{E}\rightarrow X\rightarrow\mathrm{A}\) 仍屬未經休假換班；
\(\mathrm{E}\rightarrow R\rightarrow\mathrm{A}\) 則不違反。此規則與
工作區段班型一致性不同：前者判斷同一工作區段內是否混合班型，本項則只判斷
換班前是否出現休假。

\subsubsection{\texorpdfstring{$J^M_5$}{J5}：站務配置與公平性}

以下公平性只比較目標月第一日已在職的人員。令
\[
  \mathcal E^{full}_g=\{e:g(h_e)=g,\ e\text{ 於目標月第一日已在職}\}.
\]
並令 \(\mathcal E^{full}=\bigcup_{g\in\mathcal G}\mathcal E^{full}_g\)。
對任一群組計數 $z_e$，令 $n_g=|\mathcal E^{full}_g|$、
$Z_g=\sum_{e\in\mathcal E^{full}_g}z_e$、$\bar z_g=Z_g/n_g$，並定義
\[
  F_g(z)
  =
  \sum_{e\in\mathcal E^{full}_g}
  \max\left(0,\left\lceil |z_e-\bar z_g|-\frac{3}{2}\right\rceil\right).
\]
平均值不先四捨五入。個人計數位於平均值正負 $1.5$ 的閉區間時不罰；超出後，
每遠離最近的免罰整數計數一天加 1。人數少於 2 的群組其 $F_g(z)=0$。

\paragraph{假日休假公平}
令
\[
  R^{M,hol}_{e}
  =
  \sum_{d\in\mathcal{D}^{\mathrm{tar}}\cap\mathcal{D}^{hol}}
  \rho^{M}_{ed}.
\]
則
\[
  \Phi^{M,holiday}
  =
  \sum_{g\in\mathcal{G}}F_g(R^{M,hol}).
\]

M 不計算平日休假公平。

\paragraph{非所屬站指派}
令
\[
  C^{M,other}_e
  =
  \sum_{d\in\mathcal{D}^{\mathrm{tar}}}
  \sum_{s\in\mathcal S}
  \sum_{l\in\mathcal L\setminus\{h_e\}}x^{M}_{edls}.
\]
每位人員前 8 班非所屬站正常工作免罰，總違反量為
\[
  \Phi^{M,other}
  =
  \sum_{e\in\mathcal E^M}\max(0,C^{M,other}_e-8).
\]

\paragraph{個人早小班差距}
\[
  C^{M,s}_e=\sum_{d\in\mathcal{D}^{\mathrm{tar}}}y^{M}_{eds}.
\]
對每位整月在職人員定義
\[
  \Phi^{M,EA}
  =\sum_{e\in\mathcal E^{full}}
  \max\left(0,\left|C^{M,\mathrm E}_e-C^{M,\mathrm A}_e\right|-4\right).
\]
因此早小班差 0--4 不罰，差 5 計 1，差 6 計 2，依此類推。

\paragraph{個人夜班目標}
定義夜班罰分函數
\[
  P_N(c)=
  \begin{cases}
    10,&c=0,\\
    5,&c=1,\\
    1,&c=2,\\
    0,&c\in\{3,4\},\\
    4(c-4),&c\geq5.
  \end{cases}
\]
則
\[
  \Phi^{M,N}=\sum_{e\in\mathcal E^{full}}P_N(C^{M,\mathrm N}_e).
\]
兩項班別違反量皆以個人目標計算，不再比較站群平均；\(X\) 與延伸日不計。

\subsubsection{字典序目標}

M 模型包含兩個字典序群組。第一群組衡量指定休假的滿足程度：
\[
  J^M_1=3\Phi^{M,R^*}+\Phi^{M,unusedL},
\]
第二群組將 \(J^M_4\) 的排班品質與 \(J^M_5\) 的站務配置及公平性以整數權重合併：
\[
  \begin{aligned}
    J^M_{4+5}={}&
    5\Phi^{M,ext}
    +240\Phi^{M,month,R}
    +120\Phi^{M,balance,R1}
    +20\Phi^{M,streak}\\
    &+15\Phi^{M,mixed}
    +400\Phi^{M,NE}
    +300\Phi^{M,NA}
    +5\Phi^{M,restswitch}\\
    &+\Phi^{M,other}
    +5\Phi^{M,holiday}
    +20\Phi^{M,EA}
    +50\Phi^{M,N}.
  \end{aligned}
\]
第一群組中，未滿足一日指定休假的代價為一單位未使用 \(R休\) 的三倍。第二群組的
係數則將各違反量尺度化並表達群組內交換比例；例如一單位夜--休--早型態的代價等同
八十單位外援違反量。這些交換只發生於 \(J^M_{4+5}\) 內，不得以第二群組的改善換取
\(J^M_1\) 惡化。
完整目標為
\[
  \operatorname{lexmin}
  \left(J^M_1,J^M_{4+5}\right).
\]

\section{檢修 T／YT 共用模型}

\subsection{輸入與輸出}

T 模型的輸入包括：
\begin{itemize}
  \item 目標月份與固定七日延伸區間；
  \item T 人員、專業組別、能力分數及建模日期的固定班組；
  \item 指定休假 \(R^*\)、每人目標月 \(R休\) 使用上限、固定勤務 \(X\) 與其他固定指派；
  \item 八週週期的一般 \(R\) 與 \(R1\) 額度；
  \item 依每人在職日期推導的目標月份一般 \(R\) 與 \(R1\) 基準數量；
  \item 完整前月實際班表，用以推導跨月夜轉早休假與連續狀態；
  \item 建模開始前各人員、各八週週期已使用的 R 與 R1 數量；
  \item 國定假日集合。
\end{itemize}

模型輸出每位 T 人員在目標月份內的正常出勤、\(R\)、\(R1\)、\(R休\) 或 \(X\)。正常
出勤記錄實際班別；不同於該日期固定班組時產生軟性違反量。滿足 \(R^*\) 時，仍記錄
底層實際使用的是 \(R\)、\(R1\) 或 \(R休\)。上月歷史同樣記錄逐日實際班別；延伸日期
不列入正式輸出。

\subsection{集合與索引}

\begin{description}
  \item[\(\mathcal{E}^{T}\)] T 人員集合，索引為 \(e\)。
  \item[\(\mathcal{D}^{\mathrm{tar}}\)] 目標月份內的日期集合。
  \item[\(\mathcal{D}^{+}\)] 目標月份之後固定七日的延伸日期集合。
  \item[\(\mathcal{D}\)] 完整建模日期集合，
    \(\mathcal{D}=\mathcal{D}^{\mathrm{tar}}\cup\mathcal{D}^{+}\)。
  \item[\(\mathcal{S}\)] 班別集合，
    \(\mathcal{S}=\{\mathrm{E},\mathrm{A},\mathrm{N}\}\)。
  \item[\(\mathcal{C}\)] 八週休假週期集合。
  \item[\(\mathcal{Q}\)] 非空白專業組別集合，索引為 \(q\)。
  \item[\(\mathcal{E}_{ds}\)] 日期 \(d\) 固定屬於班別 \(s\) 的人員集合。
  \item[\(\mathcal{E}^{\mathrm{tar}}_s\)] 目標月份固定屬於班別 \(s\) 的人員集合。
  \item[\(\mathcal{S}^{\mathrm{tar}}\)] 目標月份實際有人員的班別集合。
  \item[\(\mathcal{E}^{N\rightarrow E}\)] 目標月班別為早班且前月實際班表曾出現夜班的人員集合。
  \item[\(\mathcal{D}^{wd}\)] 星期一至星期五且不是國定假日的日期集合。
  \item[\(\mathcal{D}^{hol}\)] 星期六、星期日或國定假日的日期集合。
\end{description}

\subsection{參數}

令 \(\sigma_{ed}\in\mathcal{S}\) 表示人員 \(e\) 在日期 \(d\) 的固定班組。
目標月份與延伸日期都必須具有明確的 \(\sigma_{ed}\)。T 的班別時間為：早班
\(07{:}00\)--\(15{:}00\)、午班 \(15{:}00\)--\(23{:}00\)、夜班
\(23{:}00\)--翌日 \(07{:}00\)。夜班歸屬其開始日期。

令 \(\kappa_e\in\mathcal{Q}\cup\{\varnothing\}\) 為專業組別，令
\(\alpha_e\in\{1,2,3,4,5\}\) 為能力分數。依固定班組定義
\[
  \mathcal{E}_{ds}
  =
  \{e\in\mathcal{E}^{T}:\sigma_{ed}=s\},
  \qquad
  N_{ds}=|\mathcal{E}_{ds}|,
\]
以及每日正常出勤目標
\[
  B_{ds}=\left\lfloor\frac{N_{ds}}{2}\right\rfloor.
\]
班組內需要檢查的專業集合為
\[
  \mathcal{Q}_{ds}
  =
  \{q\in\mathcal{Q}:\exists e\in\mathcal{E}_{ds},\ \kappa_e=q\}.
\]

指定休假參數 \(p_{ed}\)、每人 \(R休\) 使用上限 \(k_e\)、固定勤務參數 \(z_{ed}\)、八週額度
\(G_c,H_c\)、逐人月休基準 \(n_e,m_e\) 與歷史額度
\(\overline{R}_{ec},\overline{R1}_{ec}\) 的意義與 M 模型相同。

令 \(d_0=\min\mathcal{D}^{\mathrm{tar}}-1\) 為目標月份前一日，令
\(d_1=\min\mathcal{D}^{\mathrm{tar}}\) 為目標月份第一日。對
\(e\in\mathcal{E}^{N\rightarrow E}\)，令 \(\beta_e\) 為最後一次實際夜班
結束後至 \(d_1\) 之前已出現的實際休假日數。完整前月歷史另用來
重建法規連續日數與尚未結束的實際工作區段；公平性不使用歷史累積值。

\subsection{決策變數與衍生狀態}

主要決策變數為
\[
  x^{T}_{eds}\in\{0,1\},
\]
其中 \(x^{T}_{eds}=1\) 表示人員 \(e\) 在日期 \(d\) 正常出勤班別 \(s\)。
休假變數為
\[
  r^{T}_{ed},\ r^{T,1}_{ed},\ r^{T,L}_{ed}\in\{0,1\}.
\]
定義
\[
  v_{ed}=\sum_{s\in\mathcal{S}}x^{T}_{eds},
  \qquad
  w^{T}_{ed}=v_{ed}+z_{ed},
  \qquad
  \rho^{T}_{ed}=r^{T}_{ed}+r^{T,1}_{ed}+r^{T,L}_{ed}.
\]
其中 \(X\) 是工作日，但不算正常出勤，因此不計入 T 的出勤、專業與能力指標。

令 \(q^{T,law}_{ed}\) 為只有一般 \(R\) 才歸零的法規連續日數；令
\(c^{T,work}_{ed}\) 為正常班或 \(X\) 形成的實際連續工作區段長度，\(R\) 與
\(R1\) 與 \(R休\) 都會使其歸零。兩者用途不同，不得共用同一計數。

另定義：
\begin{itemize}
  \item \(\delta^{att}_{ds}\in\mathbb{Z}_{\geq0}\)：班組出勤不足量；
  \item \(\delta^{sp}_{dsq}\in\{0,1\}\)：某專業是否完全缺席；
  \item \(\delta^{ability}_{ds}\in\{0,1,10\}\)：能力 4--5 人員不足懲罰；
  \item \(f_{ed}\in\{0,1\}\)：日期 \(d\) 是否為人員 \(e\) 在目標月份第一次正常早班；
  \item \(\delta^{bridge}_{ed}\in\mathbb{Z}_{\geq0}\)：夜班轉早班的休假不足量。
\end{itemize}

\subsection{硬性限制}

\paragraph{每日唯一指派}
\[
  \sum_{s\in\mathcal{S}}x^{T}_{eds}
  +r^{T}_{ed}+r^{T,1}_{ed}+r^{T,L}_{ed}+z_{ed}=1,
  \qquad
  \forall e\in\mathcal{E}^{T},\ d\in\mathcal{D}.
\]
固定正常班、固定 \(R\)、固定 \(R1\)、固定 \(R休\) 與固定 \(X\) 直接固定對應變數。

\paragraph{每月 R休 上限}
\[
  r^{T,L}_{ed}\leq p_{ed},
  \qquad
  \sum_{d\in\mathcal{D}^{\mathrm{tar}}}r^{T,L}_{ed}\leq k_e.
\]

\paragraph{最少十一小時休息}
令 \(\mathcal{I}^{T}\) 為所有實際休息時間少於十一小時的 T 正常班組合。對任一
\(((d,s),(d',s'))\in\mathcal{I}^{T}\)，要求
\[
  x^{T}_{eds}+x^{T}_{ed's'}\leq1,
  \qquad
  \forall e\in\mathcal{E}^{T}.
\]
正常班與固定 \(X\) 亦使用實際起訖時間檢查；兩筆固定 \(X\) 若彼此重疊或間隔
不足十一小時，輸入本身不可行。

\paragraph{連續七日至少一日一般 R}
\[
  r^{T}_{ed}=1
  \ \Longrightarrow\ 
  q^{T,law}_{ed}=0,
\]
\[
  r^{T}_{ed}=0
  \ \Longrightarrow\ 
  q^{T,law}_{ed}=q^{T,law}_{e,d-1}+1,
\]
且
\[
  0\leq q^{T,law}_{ed}\leq6.
\]
正常班、\(X\)、\(R1\) 與 \(R休\) 都會使計數增加；只有一般 \(R\) 會歸零。

\paragraph{八週休假額度}
對每個週期 \(c\)，令
\[
  \mathcal{D}_c=c\cap\mathcal{D},
  \qquad
  F_c=\left|\{d\in c:d>\max\mathcal{D}\}\right|.
\]
若模型已涵蓋週期結束日，則
\[
  \overline{R}_{ec}+\sum_{d\in\mathcal{D}_c}r^{T}_{ed}=G_c,
  \qquad
  \overline{R1}_{ec}+\sum_{d\in\mathcal{D}_c}r^{T,1}_{ed}=H_c.
\]
若尚未涵蓋週期結束日，則
\[
  \overline{R}_{ec}+\sum_{d\in\mathcal{D}_c}r^{T}_{ed}\leq G_c,
  \qquad
  \overline{R1}_{ec}+\sum_{d\in\mathcal{D}_c}r^{T,1}_{ed}\leq H_c,
\]
\[
  \overline{R}_{ec}+\sum_{d\in\mathcal{D}_c}r^{T}_{ed}+F_c\geq G_c,
\]
\[
  \overline{R1}_{ec}+\sum_{d\in\mathcal{D}_c}r^{T,1}_{ed}+F_c\geq H_c,
\]
以及
\[
  \overline{R}_{ec}+\overline{R1}_{ec}
  +\sum_{d\in\mathcal{D}_c}\left(r^{T}_{ed}+r^{T,1}_{ed}\right)
  +F_c
  \geq G_c+H_c.
\]
\(R休\) 不計入上述任一 R/R1 額度。

\paragraph{固定指派與歷史邊界}
所有固定指派均不得由模型改寫。歷史班表決定建模區間第一日前的
\(q^{T,law}\)、\(c^{T,work}\)、八週累積值與跨月夜轉早所需狀態。

\subsection{軟性目標與權重}

\subsubsection{\texorpdfstring{$J^T_1$}{J1}：指定休假}

\paragraph{指定休假滿足}
\[
  \Phi^{T,R^*}
  =
  \sum_{\substack{e\in\mathcal{E}^{T},\ d\in\mathcal{D}^{\mathrm{tar}}\\p_{ed}=1}}
  \left(1-r^{T}_{ed}-r^{T,1}_{ed}-r^{T,L}_{ed}\right).
\]

\paragraph{未使用 R休 額度}
未使用的 \(R休\) 上限違反量為
\[
  \Phi^{T,unusedL}
  =\sum_{e\in\mathcal{E}^{T}}
  \left(k_e-\sum_{d\in\mathcal{D}^{\mathrm{tar}}}r^{T,L}_{ed}\right).
\]

\subsubsection{\texorpdfstring{$J^T_2$}{J2}：班組人力品質}

\paragraph{月班別一致性}
每個模型日期排入非該日期固定班組的正常工作格計一次違反；\(X\) 不計。令
\[
  \Phi^{T,nonmonthly}
  =
  \sum_{e\in\mathcal{E}^{T}}
  \sum_{d\in\mathcal{D}}
  \sum_{\substack{s\in\mathcal{S}\\s\neq\sigma_{ed}}}x^{T}_{eds}.
\]
目標月以輸入的月班組為基準；延伸日期以早 \(\rightarrow\) 午 \(\rightarrow\) 夜
\(\rightarrow\) 早輪轉後的下月班組為基準。

\paragraph{班組出勤人數}
日期 \(d\)、班別 \(s\) 的正常出勤人數為
\[
  A_{ds}
  =
  \sum_{e\in\mathcal{E}^{T}}x^{T}_{eds}.
\]
要求
\[
  \delta^{att}_{ds}\geq B_{ds}-A_{ds},
  \qquad
  \delta^{att}_{ds}\geq0.
\]
總違反量為
\[
  \Phi^{T,att}
  =
  \sum_{d\in\mathcal{D}^{\mathrm{tar}}}
  \sum_{s\in\mathcal{S}}\delta^{att}_{ds}.
\]

\paragraph{專業人員出勤}
對每個 \(q\in\mathcal{Q}_{ds}\)，令
\[
  A^{sp}_{dsq}
  =
  \sum_{\substack{e\in\mathcal{E}^{T}\\\kappa_e=q}}x^{T}_{eds},
\]
並令
\[
  M^{sp}_{dsq}
  =
  \left|\{e\in\mathcal{E}^{T}:\kappa_e=q\}\right|.
\]
以 \(\delta^{sp}_{dsq}=1\) 表示該專業無人正常出勤：
\[
  A^{sp}_{dsq}\geq1-\delta^{sp}_{dsq},
  \qquad
  A^{sp}_{dsq}\leq M^{sp}_{dsq}(1-\delta^{sp}_{dsq}).
\]
總違反量為
\[
  \Phi^{T,specialty}
  =
  \sum_{d\in\mathcal{D}^{\mathrm{tar}}}
  \sum_{s\in\mathcal{S}}
  \sum_{q\in\mathcal{Q}_{ds}}\delta^{sp}_{dsq}.
\]
專業空白的人員仍可正常出勤，但空白不形成需要覆蓋的專業類別。

\paragraph{高能力人員配置}
日期 \(d\)、班別 \(s\) 中，能力至少為 4 的正常出勤人數為
\[
  H_{ds}
  =
  \sum_{\substack{e\in\mathcal{E}^{T}\\\alpha_e\geq4}}x^{T}_{eds}.
\]
每班希望至少有兩位能力 4--5 的正常出勤者，其不足懲罰為
\[
  \delta^{ability}_{ds}
  =
  \begin{cases}
    10, & H_{ds}=0,\\
    1, & H_{ds}=1,\\
    0, & H_{ds}\geq2.
  \end{cases}
\]
總違反量為
\[
  \Phi^{T,ability}
  =
  \sum_{d\in\mathcal{D}^{\mathrm{tar}}}
  \sum_{s\in\mathcal{S}}\delta^{ability}_{ds}.
\]
能力 4--5 的跨班人員按實際工作班別計入。

\subsubsection{\texorpdfstring{$J^T_3$}{J3}：休假分布}

\paragraph{每月一般 R 分布}
令
\[
  Q^{T,R}_e=\sum_{d\in\mathcal{D}^{\mathrm{tar}}}r^{T}_{ed}.
\]
沿用
\[
  f_t(q)=|q-t|^2,
\]
則
\[
  \Phi^{T,month,R}
  =
  \sum_{e\in\mathcal{E}^{T}}f_{n_e}\!\left(Q^{T,R}_e\right),
\]

\paragraph{八週累積 R1 餘額}
令 \(p^{T,1}_{eI}\)、\(a_{eI}\) 與 \(T^{T,1}_{eI}\) 使用和 M 相同的跨區間
比例分配與新進折抵定義，則
\[
  A^{T,1}_{eI}=p^{T,1}_{eI}
  +\sum_{d\in\mathcal D^{\mathrm{tar}}\cap I}r^{T,1}_{ed},
  \qquad
  \Phi^{T,balance,R1}=\sum_{e\in\mathcal E^T}\sum_I
  \left(A^{T,1}_{eI}-T^{T,1}_{eI}\right)^2.
\]
此處採對稱平方偏差，亦即提前使用與尚未使用的每一單位餘額皆計入懲罰；因此在硬性
限制可行的範圍內，累積使用量會趨近分配目標。

\subsubsection{\texorpdfstring{$J^T_4$}{J4}：工作型態品質}

\paragraph{連續工作區段長度}
T 的實際工作區段同樣由連續正常班或 \(X\) 組成，\(R\)、\(R1\) 與 \(R休\)
都會結束區段。其長度懲罰為：
\[
  D^T(L)=
  \begin{cases}
    4, & L=1,\\
    1, & L=2,\\
    0, & L\in\{3,4\},\\
    1, & L=5,\\
    2(L-4), & L\geq6.
  \end{cases}
\]
令 \(\mathcal{B}^{T,work}_e\) 為承接歷史後的
最大實際工作區段集合，則
\[
  \Phi^{T,streak}
  =
  \sum_{e\in\mathcal{E}^{T}}
  \sum_{B\in\mathcal{B}^{T,work}_e:\ \max B\in\mathcal{D}^{\mathrm{tar}}}
  D^T(|B|).
\]
此違反量控制合理工作區段長度；\(R休\) 與 \(R1\) 同樣會切斷工作區段，但不重置
只有一般 \(R\) 才歸零的法規連續日數。

\paragraph{跨月夜轉早休假}
此規則只適用於 \(e\in\mathcal{E}^{N\rightarrow E}\)。令 \(f_{ed}=1\) 表示日期
\(d\) 是該人在目標月份第一次正常早班。若 \(f_{ed}=1\)，最後夜班至該日之前
的休假日數為
\[
  B^{bridge}_{ed}
  =
  \beta_e
  +
  \sum_{\substack{\tau\in\mathcal{D}^{\mathrm{tar}}\\\tau<d}}
  \rho^{T}_{e\tau}.
\]
定義
\[
  \delta^{bridge}_{ed}
  =
  \begin{cases}
    \max(0,2-B^{bridge}_{ed}), & f_{ed}=1,\\
    0, & f_{ed}=0.
  \end{cases}
\]
總違反量為
\[
  \Phi^{T,monthrest}
  =
  \sum_{e\in\mathcal{E}^{N\rightarrow E}}
  \sum_{d\in\mathcal{D}^{\mathrm{tar}}}\delta^{bridge}_{ed}.
\]
此規則要求夜班轉早班的人員，在最後夜班與第一次早班之間至少有兩個實際休假日。

\paragraph{月交界休假平衡}
令 \(\overline{\rho}^{T}_{e,d_0}\) 為前月最後一日的歷史休假指示，並定義
\[
  B^{-}
  =
  \sum_{e\in\mathcal{E}^{N\rightarrow E}}
  \overline{\rho}^{T}_{e,d_0},
  \qquad
  B^{+}
  =
  \sum_{e\in\mathcal{E}^{N\rightarrow E}}
  \rho^{T}_{e,d_1}.
\]
月交界休假分散違反量為
\[
  \Phi^{T,monthbalance}=|B^{-}-B^{+}|.
\]

\subsubsection{\texorpdfstring{$J^T_5$}{J5}：休假公平}

\paragraph{平日休假公平}
令
\[
  R^{T,wd}_{e}
  =
  \sum_{d\in\mathcal{D}^{\mathrm{tar}}\cap\mathcal{D}^{wd}}
  \rho^{T}_{ed}.
\]
只在目標月份同一固定班組內比較：
\[
  \Phi^{T,weekday}
  =
  \sum_{s\in\mathcal{S}^{\mathrm{tar}}}
  \left(
    \max_{e\in\mathcal{E}^{\mathrm{tar}}_s}R^{T,wd}_{e}
    -
    \min_{e\in\mathcal{E}^{\mathrm{tar}}_s}R^{T,wd}_{e}
  \right).
\]

\paragraph{假日休假公平}
令
\[
  R^{T,hol}_{e}
  =
  \sum_{d\in\mathcal{D}^{\mathrm{tar}}\cap\mathcal{D}^{hol}}
  \rho^{T}_{ed}.
\]
對每個月班別 $s$，令 $n_s=|\mathcal E^{\mathrm{tar}}_s|$、
$A_s=\sum_{e\in\mathcal E^{\mathrm{tar}}_s}R^{T,hol}_e/n_s$。則
\[
  \Phi^{T,holiday}
  =
  \sum_{s\in\mathcal{S}^{\mathrm{tar}}}
  \sum_{e\in\mathcal E^{\mathrm{tar}}_s}
  \max\left(0,
    \left\lceil |R^{T,hol}_e-A_s|-\frac{3}{2}\right\rceil
  \right).
\]
平均值不先四捨五入；個人計數位於平均值正負 $1.5$ 天內不罰，超出後按天線性增加。

\subsubsection{字典序目標}

T 模型定義五個依序最佳化的目標群組：
\[
  J^T_1=3\Phi^{T,R^*}+\Phi^{T,unusedL},
\]
\[
  J^T_2
  =
  9\Phi^{T,nonmonthly}
  +9\Phi^{T,att}
  +3\Phi^{T,specialty}
  +\Phi^{T,ability},
\]
\[
  J^T_3
  =
  \Phi^{T,month,R}
  +\Phi^{T,balance,R1},
\]
\[
  J^T_4
  =
  3\Phi^{T,streak}
  +12\Phi^{T,monthrest}
  +5\Phi^{T,monthbalance},
\]
\[
  J^T_5
  =
  2\Phi^{T,weekday}
  +4\Phi^{T,holiday}.
\]
完整目標為
\[
  \operatorname{lexmin}
  \left(J^T_1,J^T_2,J^T_3,J^T_4,J^T_5\right).
\]

\section{求解流程與候選方案}

\subsection{M 萬年班表提示}

對站務模型，可給定長度為 56 日的週期模板 \(\mathbf h\) 作為初始搜尋提示。模板起點
與每個八週週期的首日對齊；空白位置不指定值，固定指派及指定休假位置亦不採用模板值。
令 \(\mathcal H\) 為可提示的變數索引集合，則提示可寫為
\[
  x_i\leftarrow h_i,
  \qquad i\in\mathcal H.
\]
此關係不是 \(\mathcal F\) 的限制，僅引導 CP-SAT 優先探索與模板接近的分支。求得第一個
可行解 \(\mathbf x^{(0)}\) 後，模板提示由完整 incumbent 提示
\(x_i\leftarrow x_i^{(0)}\) 取代，再進行各字典序階段。因此模板即使與歷史狀態或
最佳解不一致，也不會縮小可行域或改變最優解集合；其作用僅限於改善初始可行解搜尋。

\subsection{差異候選}

每次求解最多保留三份候選。令 \(\mathcal U\) 為目標月份中未被固定指派的逐人逐日
格集合，\(s_i(\mathbf x)\) 為候選 \(\mathbf x\) 在格 \(i\) 的實際狀態。固定工作、固定
休假、固定勤務與延伸日期不屬於 \(\mathcal U\)；指定休假仍須由模型選擇底層休假種類，
故屬於比較集合。對兩候選 \(\mathbf x\) 與 \(\mathbf y\)，定義離散漢明距離
\[
  \Delta(\mathbf x,\mathbf y)
  =\sum_{i\in\mathcal U}\mathbf 1\!\left[s_i(\mathbf x)\neq s_i(\mathbf y)\right].
\]
申請標記不構成額外狀態，但 \(R\)、\(R1\) 與 \(R休\) 彼此不同。

令 \(N_{cell}=|\mathcal U|\)，最低差異格數為
\[
  \Delta_{min}=\left\lceil0.05N_{cell}\right\rceil.
\]
若已取得候選集合 \(\mathcal C\)，新候選 \(\mathbf x\) 必須滿足
\[
  \Delta(\mathbf x,\mathbf y)\geq\Delta_{min},
  \qquad \forall\mathbf y\in\mathcal C.
\]
替代候選搜尋使用主要最佳化完成後的剩餘時間。首先固定所有具名目標分數，使新候選
與主候選同分。T 僅在所有目標群組均已證明最優時執行此步。M 若未能於配置時間內找到
第二份同分候選，可將各具名目標分別放寬至主候選的 120\%。若主候選第 \(j\) 項分數為
\(V_j^{(1)}\)，候選 \(k\) 須滿足
\[
  5V_j^{(k)}\leq6V_j^{(1)},\qquad\forall j.
\]
此整數式避免百分比取整歧義；當 \(V_j^{(1)}=0\) 時亦強制 \(V_j^{(k)}=0\)。因此候選
多樣性由成對距離保證，品質則由同分等式或逐項上界控制，而非另行改變原目標權重。

\section{數學模型與 CP-SAT 實作對照}

數學式必須轉換為 CP-SAT 可接受的有限整數表示。表中整理主要數學物件及其編碼方式；
每一列皆保持等價關係，而非以連續鬆弛近似原條件。

\footnotesize
\setlength{\tabcolsep}{3pt}
\begin{longtable}{p{3.6cm}p{5.0cm}p{5.8cm}}
  \toprule
  數學物件 & CP-SAT 表示 & 等價性說明 \\
  \midrule
  \endfirsthead
  \toprule
  數學物件 & CP-SAT 表示 & 等價性說明 \\
  \midrule
  \endhead
  每日唯一狀態 & 布林變數之和等於 1 & 恰一個工作、休假或固定勤務狀態取值為 1 \\
  固定指派 & 對應布林變數等於 1 & 配合每日唯一狀態，自動排除其餘狀態 \\
  狀態轉移 & 前一日整數狀態加條件式等式 & 由休假布林值選擇歸零或遞增分支 \\
  \(\max(0,z)\) & \texttt{MaxEquality} & 輔助變數精確等於 \(z\) 與 0 的最大值 \\
  \(|z|\) & \texttt{MaxEquality} & 輔助變數精確等於 \(z\) 與 \(-z\) 的最大值 \\
  分段懲罰 \(D(L)\) & \texttt{Element} 查表限制 & 對有限整數 \(L\) 直接選取唯一函數值 \\
  模式指示變數 & 重述的合取限制 & 指示變數取 1 若且唯若模式中各狀態同時成立 \\
  群組最大值與最小值 & \texttt{MaxEquality}、\texttt{MinEquality} & 精確取得群組成員計數的極值 \\
  字典序目標 & 逐階段最小化及 \(J_p=J_p^*\) & 完成階段的最優值成為後續階段硬限制 \\
  候選漢明距離 & 狀態相異指示變數之和 & 精確計算未固定日格的不同數量 \\
  \bottomrule
\end{longtable}
\normalsize
\setlength{\tabcolsep}{6pt}

\subsection{實際求解程序}

對每個隨機種子，求解程序依下列順序使用同一份總時間預算：
\begin{enumerate}
  \item 驗證集合、索引、時間區間與歷史狀態的完備性，建立目標月及七日延伸區間。
  \item 建立決策變數、衍生狀態與全部硬性限制，先搜尋任一 \(\mathbf x^{(0)}\in\mathcal F\)。
  \item 以 \(\mathbf x^{(0)}\) 作為 incumbent 提示，建立各違反量及目標群組。
  \item 對 \(p=1,\ldots,P\) 依序最小化 \(J_p\)；僅在第 \(p\) 階段取得最優性證明時，
    加入 \(J_p=J_p^*\) 並繼續下一階段。
  \item 以剩餘時間加入成對漢明距離限制，搜尋至多三份候選。
\end{enumerate}

令每個種子的時間上限為 \(\tau\)，搜尋 worker 數為 \(W\)。站務模型可依序執行
\(K_M\geq1\) 個隨機種子，檢修模型取 \(K_T=1\)。當 \(K_M>1\) 時，不同種子的結果
以第一候選的目標向量進行字典序比較並選取較小者；若目標向量相同，則不因種子編號
改變品質判定。此 portfolio 策略利用 CP-SAT 搜尋樹對分支順序的敏感性，提高有限時間內
取得良好 incumbent 的機率，但 \((\tau,W,K_M)\) 均屬求解參數，不改變模型本身。

\subsection{狀態解讀與限制}

\begin{description}
  \item[\texttt{Optimal}] 每個字典序階段均已證明最優；替代候選搜尋是否耗盡時間不影響此證明。
  \item[\texttt{TimeLimit}] 總時間預算耗盡。若附帶 incumbent，該班表滿足全部硬性限制，
    但尚未完成階段的目標值只代表已知上界，不構成最優性證明。
  \item[\texttt{Infeasible}] 可行域 \(\mathcal F=\varnothing\)，亦即硬性限制無法同時成立。
  \item[\texttt{InvalidInput}] 集合、參數或歷史邊界不滿足模型定義域，故不建立最佳化問題。
\end{description}

對最小化階段，若 CP-SAT 同時提供 incumbent 值 \(U_p\) 與最佳界 \(L_p\)，則
\(L_p\leq J_p^*\leq U_p\)。因此有限時間結果除列示違反量外，亦可用
\(U_p-L_p\) 判讀尚未證明的絕對最佳化缺口；只有 \(U_p=L_p\) 時該階段方告完成。

\section{結論}

本文將站務與檢修月排班表述為有限域上的字典序限制最佳化問題。硬性限制刻畫法規、
時間邊界、額度與營運可行性，非負違反量則量化指定休假、工作型態、人力品質及群組
公平。M／YM 與 T／YT 分別採用符合其作業特性的模型，同時共享一致的休假語意與跨月
處理原則。七日延伸區間及完整歷史狀態避免月界截斷造成的假可行解；逐階段固定最優值
則嚴格維持管理目標的優先次序。

CP-SAT 編碼利用布林重述、最大值等式與有限域查表精確表示非線性離散邏輯，不需連續
近似。週期模板與多種子 portfolio 改善有限時間內的搜尋效率而不改變可行域；候選生成
再以成對漢明距離與逐項品質界提供具有實質差異的替代方案。由於每個目標分數皆可還原
為違反量與權重，所得班表不僅可驗證其合法性，也能清楚說明不同可行方案之間的品質取捨。

\end{document}
